% source: J. S. Milne, "Abelian Varieties" (course notes), v2.00, 2008, www.jmilne.org
% pdf_page: 0025
% doc_page: 19 (Chapter I, Section 3 -- proof of Theorem 3.4; Corollary 3.7; Theorem 3.8; Proposition 3.9 and beginning of its proof)
% retrieved: 2026-09-04
% transcribed_by: codex-gpt-5 (vision, rendered at 200dpi via pdftoppm; checked at 350dpi)

% [running head: 3. RATIONAL MAPS INTO ABELIAN VARIETIES, page 19]

\DeclareMathOperator{\Image}{Im}

% [continuation of Theorem 3.4 from the previous page]

\paragraph{Proof (of 3.4).}
Clearly we can assume $k$ to be algebraically closed. Consider first the case
that $V$ has dimension 1. From the (3.5), we know that $V$ can be embedded into
a nonsingular complete curve $C$, and (3.2) shows that $\alpha$ extends to a map
$\bar{\alpha}\colon C \times W \to A$. Now the Rigidity Theorem (1.1) shows that
$\bar{\alpha}$ is constant. In the general case, let $C$ be an irreducible curve
on $V$ passing through $v_0$ and nonsingular at $v_0$, and let $C' \to C$ be the
normalization of $C$. By composition, $\alpha$ defines a morphism
$C' \times W \to A$, which the preceding argument shows to be constant.
Therefore $\alpha(C \times W) = \{a_0\}$, and Lemma 3.6 completes the proof.
\hfill $\square$

\paragraph{Corollary 3.7.}
\textit{Every rational map $\alpha\colon G \dashrightarrow A$ from a group
variety to an abelian variety is the composite of a homomorphism
$h\colon G \to A$ with a translation.}

\paragraph{Proof.}
Theorem 3.2 shows that $\alpha$ is a regular map. The rest of the proof is the
same as that of Corollary 1.2. \hfill $\square$

\subsection*{Abelian varieties up to birational equivalence.}

A rational map $\varphi\colon V \dashrightarrow W$ is \textit{dominating} if
$\Image(\varphi_U)$ is dense in $W$ for one (hence all) representatives
$(U, \varphi_U)$ of $\varphi$. Then $\varphi$ defines a homomorphism
$k(W) \to k(V)$, and every such homomorphism arises from a (unique) dominating
rational map (exercise!).

A rational map $\varphi$ is \textit{birational} if the corresponding
homomorphism $k(W) \to k(V)$ is an isomorphism. Equivalently, if there exists a
rational map $\psi\colon W \dashrightarrow V$ such that $\varphi \circ \psi$
and $\psi \circ \varphi$ are both the identity map wherever they are defined.
Two varieties $V$ and $W$ are \textit{birationally equivalent} if there exists
a birational map $V \dashrightarrow W$; equivalently, if $k(V) \approx k(W)$.

In general, two varieties can be birationally equivalent without being
isomorphic (see the start of this section 1 for examples). In fact, every
variety (even complete and nonsingular) of dimension $> 1$ will be birationally
equivalent to many nonisomorphic varieties. However, Theorem 3.1 shows that two
complete nonsingular curves that are birationally equivalent will be isomorphic.
The same is true of abelian varieties.

\paragraph{Theorem 3.8.}
\textit{If two abelian varieties are birationally equivalent, then they are
isomorphic (as abelian varieties).}

\paragraph{Proof.}
Let $A$ and $B$ be the abelian varieties. A rational map
$\varphi\colon A \dashrightarrow B$ extends to a regular map $A \to B$ (by
3.2). If $\varphi$ is birational, its inverse $\psi$ also extends to a regular
map, and the composites $\varphi \circ \psi$ and $\psi \circ \varphi$ will be
identity maps because they are on open sets. Hence there is an isomorphism
$\alpha\colon A \to B$ of algebraic varieties. After composing it with a
translation, it will map 0 to 0, and then Corollary 1.2 shows that it preserves
the group structure. \hfill $\square$

\paragraph{Proposition 3.9.}
\textit{Every rational map $\mathbb{A}^1 \dashrightarrow A$ or
$\mathbb{P}^1 \dashrightarrow A$ is constant.}

\paragraph{Proof.}
According to (3.2), $\alpha$ extends to a regular map on the whole of
$\mathbb{A}^1$. After composing $\alpha$ with a translation, we may suppose
that $\alpha(0) = 0$. Then $\alpha$ is a homomorphism,
\[
  \alpha(x + y) = \alpha(x) + \alpha(y), \qquad
  \text{all } x,y \in \mathbb{A}^1(k^{\mathrm{al}}) = k^{\mathrm{al}}.
\]

% [proof of Proposition 3.9 continues on the next PDF page]
